\documentclass[11pt]{article}
\usepackage[margin=1in]{geometry}
\usepackage{enumitem}
\usepackage{hyperref}
\usepackage{listings}
\usepackage{xcolor}
\usepackage{booktabs}
\usepackage{titlesec}
\usepackage{fancyhdr}

\lstset{
 basicstyle=\ttfamily\small,
 keywordstyle=\color{blue},
 commentstyle=\color{gray},
 stringstyle=\color{red},
 breaklines=true,
 columns=flexible,
 literate={_}{{\_}}1
}

\titleformat{\section}{\normalfont\Large\bfseries}{}{0em}{}
\titleformat{\subsection}{\normalfont\large\bfseries}{}{0em}{}
\titleformat{\subsubsection}{\normalfont\normalsize\bfseries\itshape}{}{0em}{}

\pagestyle{fancy}
\fancyhf{}
\fancyhead[L]{\small\textbf{Software Guide}}
\fancyhead[R]{\small\textit{Data Structures, Algorithms, and Applications in C++}}
\fancyfoot[C]{\thepage}

\setlength{\parindent}{0pt}
\setlength{\parskip}{4pt}

\begin{document}

\begin{center}
{\LARGE\bfseries Software Guide}\\[6pt]
{\large Data Structures, Algorithms, and Applications in C++}\\[4pt]
{\small 76 header-only implementations $\cdot$ dsa namespace $\cdot$ C++23}
\end{center}

\vspace{6pt}
\noindent\hrulefill
\vspace{6pt}

Examples show the public API only; include the header named in each entry (and any standard-library headers used by the example) before compiling them.

%% ============================================================
%% PART II: DATA STRUCTURES
%% ============================================================

\section{Part II --- Data Structures}

% ---- LINEAR LISTS ----
\subsection{Linear Lists}

\subsubsection{Array List}
\textbf{Header:} code/include/array\_list.h\\
\textbf{Class:} dsa::array\_list$<$T$>$\\

Dynamic array with amortized append. Grows via reallocation. This is the most fundamental sequential container, analogous to std::vector but built from scratch for pedagogical clarity.

\textbf{Why:} You need fast random access and appending, and don't need insertion/deletion in the middle. \par \textbf{Where to use:} Use when building other data structures (stacks, queues, deques), managing collections with index access, or as scratch buffers. Default choice for sequential data.

\textbf{Input:} Elements of type T (must satisfy std::regular)\\
\textbf{Output:} Index-based access via operator[]; iterators for traversal

\begin{lstlisting}
dsa::array_list<int> a;
a.push_back(10);
a.push_back(20);
a[0]; // 10
a.size(); // 2
\end{lstlisting}

\subsubsection{Linked List}
\textbf{Header:} code/include/linked\_list.h\\
\textbf{Class:} dsa::linked\_list$<$T$>$\\

Singly linked list using std::unique\_ptr for automatic memory management. Supports forward iteration and an in-place \texttt{reverse()} operation.

\textbf{Why:} You need insert/delete at known positions without shifting elements, and don't need random access. \par \textbf{Where to use:} Use for adjacency lists in sparse graphs, undo/redo buffers, implementing other containers, and when elements are frequently inserted/deleted in the middle.

\textbf{Input:} Elements of any copyable/movable type\\
\textbf{Output:} insert at front; iteration via range-for

\begin{lstlisting}
dsa::linked_list<int> l;
l.push_front(3);
l.push_back(4);
for (auto x : l) { /* 3, 4 */ }
\end{lstlisting}

\subsubsection{Doubly Linked List}
\textbf{Header:} code/include/doubly\_linked\_list.h\\
\textbf{Class:} dsa::doubly\_linked\_list$<$T$>$\\

Doubly linked list with sentinel node for clean boundary handling. Bidirectional traversal.

\textbf{Why:} You need insert/delete at both ends AND delete at arbitrary positions (given iterator). \par \textbf{Where to use:} Use for LRU caches (with hash map), browser history, text editor buffers, and deques. Better than singly linked list when you need backward traversal.

\begin{lstlisting}
dsa::doubly_linked_list<int> dll;
dll.push_back(1);
dll.push_back(2);
dll.push_front(0);
dll.front(); // 0
dll.back(); // 2
dll.pop_back();
for (auto x : dll) { /* 0, 1 */ }
\end{lstlisting}

\subsubsection{Bit Vector}
\textbf{Header:} code/include/bit\_vector.h\\
\textbf{Class:} dsa::BitVector\\

Compact bit storage with bulk operations. Uses native word-level operations for counting set bits.

\textbf{Why:} You need compact storage for boolean/bit flags with fast bulk operations (popcount, find-first). \par \textbf{Where to use:} Use for visited arrays in BFS/DFS, bloom filter internals, database bitmap indexes, permission flags, and any compact boolean set over a small universe.

\begin{lstlisting}
dsa::BitVector bv(8);
bv.set(3); bv.set(5);
bv.count(); // 2
bv.find_first(); // 3
bv.flip(3);
bv.get(3); // false
\end{lstlisting}

% ---- SPARSE ----
\subsubsection{Sparse Matrix}
\textbf{Header:} code/include/sparse\_matrix.h\\
\textbf{Class:} dsa::SparseMatrix$<$T$>$\\

Sparse matrix in COO (coordinate) and CSR (compressed sparse row) format. Supports transpose and matrix-vector multiply.

\textbf{Why:} Your matrix is mostly zeros and storing it densely wastes memory and compute. \par \textbf{Where to use:} Use for large-scale scientific computing, ML feature matrices, graph adjacency matrices, and FEM/PDE solvers where nnz << rows*cols.

\begin{lstlisting}
dsa::SparseMatrix<double> m(3, 3);
m.set(0, 0, 1.0);
m.set(1, 2, 3.0);
auto t = m.transpose();
std::vector<double> x = {1, 2, 3};
auto y = m.multiply(x); // matrix-vector product
\end{lstlisting}

\subsubsection{Sparse Vector}
\textbf{Header:} code/include/sparse\_vector.h\\
\textbf{Class:} dsa::SparseVector$<$T$>$\\

\textbf{Why:} Your vector is high-dimensional but mostly zeros, and dense storage wastes memory. \par \textbf{Where to use:} Use for TF-IDF vectors in IR, ML feature representations, recommendation system embeddings, and any sparse high-dimensional data.

\begin{lstlisting}
dsa::SparseVector<double> sv(100);
sv.set(5, 2.0);
sv.set(99, 3.0);
double d = sv.dot(sv); // 13.0
auto s = sv.scale(2.0);
\end{lstlisting}

% ---- STACKS ----
\subsection{Stacks}
\textbf{Header:} code/include/stacks.h\\
\textbf{Classes:} dsa::array\_stack$<$T$>$, dsa::linked\_stack$<$T$>$\\

Two implementations: vector-backed array stack and pointer-based linked stack. LIFO access pattern.

\textbf{Why:} You need LIFO access: last-in, first-out. \par \textbf{Where to use:} Use for expression evaluation, DFS traversal, backtracking, undo systems, parenthesis matching, and converting infix to postfix notation.

\begin{lstlisting}
dsa::array_stack<int> s;
s.push(1);
s.push(2);
s.top(); // 2
s.pop();
s.top(); // 1
\end{lstlisting}

% ---- QUEUES ----
\subsection{Queues \& Deques}
\textbf{Header:} code/include/queues.h\\
\textbf{Classes:} dsa::circular\_queue$<$T$>$, dsa::linked\_queue$<$T$>$\\

FIFO access pattern. Circular queue uses a dynamically growing ring buffer while preserving FIFO order.

\textbf{Why:} You need FIFO access: first-in, first-out. \par \textbf{Where to use:} Use for BFS traversal, task scheduling, print spooling, producer-consumer patterns, and request handling in servers.

\subsubsection{Deque}
\textbf{Header:} code/include/deque.h\\
\textbf{Class:} dsa::Deque$<$T$>$\\

Double-ended queue (block-based circular buffer). Combines the benefits of arrays and linked lists.

\textbf{Why:} You need push/pop at both ends AND random access by index. \par \textbf{Where to use:} Use for sliding window algorithms, BFS with deques, implementing std::deque, and when you need both stack and queue behavior.

\begin{lstlisting}
dsa::Deque<int> dq;
dq.push_back(1);
dq.push_back(2);
dq.push_front(0);
dq.at(0); // 0
dq.size(); // 3
dq.pop_front();
\end{lstlisting}

% ---- HASHING ----
\subsection{Hashing}

\subsubsection{Hash Table (Chaining)}
\textbf{Header:} code/include/hash\_table.h\\
\textbf{Class:} dsa::hash\_table\_chaining$<$K,V$>$\\

Hash table with separate chaining. Each bucket is a linked list handling collisions gracefully. Supports iterators for traversal.

\textbf{Why:} You need average-case insert/find/erase by key, and don't need ordering. \par \textbf{Where to use:} Use for symbol tables, caches, frequency counting, database indexes, and any key-value lookup where order doesn't matter.

\begin{lstlisting}
dsa::hash_table_chaining<std::string,int> h;
h.insert("one", 1);
h.find("one"); // 1
h.erase("one");
\end{lstlisting}

\subsubsection{Open Addressing Hash}
\textbf{Header:} code/include/open\_addressing\_hash.h\\
\textbf{Classes:} linear probing, quadratic probing, double hashing\\

All elements stored in a single array. Three probing strategies for collision resolution. Better cache performance than chaining when load factor is moderate.

\textbf{Why:} You need a hash table with better cache locality than chaining, and can tolerate slightly higher collision costs. \par \textbf{Where to use:} Use when memory is constrained (no pointer overhead), for cache-friendly hash maps, and in real-time systems. Linear probing is simplest; double hashing gives best distribution.

\begin{lstlisting}
dsa::open_addressing_hash<std::string,int> h;
h.insert("one", 1);
h.insert("two", 2);
auto* val = h.find("one"); // points to 1
h.erase("two");
h.contains("two"); // false
\end{lstlisting}

\subsubsection{Skip List}
\textbf{Header:} code/include/skip\_list.h\\
\textbf{Class:} dsa::skip\_list$<$T$>$\\

Randomized ordered structure providing balanced BST performance without rotations. Multiple levels of forward pointers give expected time.

\textbf{Why:} You need ordered operations with simple implementation, and can accept probabilistic balance. \par \textbf{Where to use:} Use as a simpler alternative to balanced BSTs when exact balance isn't critical. This implementation is single-threaded; use a separately synchronized or concurrent skip list when concurrent access is required.

\begin{lstlisting}
dsa::skip_list<int> sl;
sl.insert(5);
sl.insert(3);
sl.insert(7);
sl.contains(5); // true
sl.erase(3);
auto v = sl.to_sorted_vector(); // [5, 7]
\end{lstlisting}

% ---- TREES ----
\subsection{Trees}

\subsubsection{Binary Tree}
\textbf{Header:} code/include/binary\_tree.h\\
\textbf{Class:} dsa::binary\_tree$<$T$>$\\

Binary tree with four traversals (inorder, preorder, postorder, level-order), Morris traversal (\(O(1)\) auxiliary space), and BST insert/erase/find.

\textbf{Why:} You need a general tree structure for hierarchical data or as a foundation for other algorithms. \par \textbf{Where to use:} Use as the building block for BSTs, heaps, Huffman trees, expression parsing, and decision trees. Start here before moving to balanced variants.

\begin{lstlisting}
dsa::binary_tree<int> bt;
bt.insert(5);
bt.insert(3);
bt.insert(7);
bt.contains(3); // true
bt.erase(3);
\end{lstlisting}

\subsubsection{Expression Tree}
\textbf{Header:} code/include/expression\_tree.h\\
\textbf{Functions:} dsa::build\_expr\_tree, dsa::eval\_expr\_tree\\

Build from postfix notation, evaluate arithmetic expressions, convert between infix/postfix/preorder notations.

\textbf{Why:} You need to parse, evaluate, or transform arithmetic expressions. \par \textbf{Where to use:} Use in compilers, calculators, symbolic math systems, and query optimizers for expression evaluation and notation conversion.

\begin{lstlisting}
std::vector<std::string> postfix = {"3", "4", "+", "2", "*"};
auto* root = dsa::build_expr_tree(postfix);
double val = dsa::eval_expr_tree(root); // 14.0
\end{lstlisting}

% ---- HEAPS ----
\subsection{Heaps \& Priority Queues}

\subsubsection{Max Heap}
\textbf{Header:} code/include/max\_heap.h\\
\textbf{Class:} dsa::max\_heap$<$T$>$\\

Standard binary max-heap implemented as an array. Also provides heap\_sort() for in-place sorting.

\textbf{Why:} You need fast access to the maximum element with insert/delete. \par \textbf{Where to use:} Use as the default priority queue, for heap sort, finding k-th largest, scheduling, and Dijkstra's algorithm. Also useful for median maintenance with two heaps.

\begin{lstlisting}
dsa::max_heap<int> h;
h.push(3); h.push(1); h.push(4);
h.top(); // 4
h.pop();
h.top(); // 3
std::vector<int> v = {5, 1, 3, 2};
dsa::heap_sort(v); // v = {1, 2, 3, 5}
\end{lstlisting}

\subsubsection{D-ary Heap}
\textbf{Header:} code/include/dary\_heap.h\\
\textbf{Class:} dsa::dary\_heap$<$T,D$>$\\

Generalized d-children min-heap (default D=4). A higher branching factor can reduce heap height, but this implementation provides push/pop/top only; it does not implement decrease-key.

\textbf{Why:} You need a simple d-ary priority queue with push/pop/top operations. \par \textbf{Where to use:} Use as a teaching implementation or when a reduced heap height is useful. For indexed priority updates, use the indexed priority queue instead.

\begin{lstlisting}
dsa::dary_heap<int, 4> h;
h.push(3); h.push(1); h.push(4);
h.top(); // 1 (min-heap)
h.pop();
\end{lstlisting}

\subsubsection{Indexed Priority Queue}
\textbf{Header:} code/include/indexed\_pq.h\\
\textbf{Classes:} dsa::indexed\_min\_pq$<$T$>$, dsa::indexed\_max\_pq$<$T$>$\\

Priority queue that supports updating priorities by index. Tracks which element is at each position.

\textbf{Why:} You need a priority queue where you can update the priority of any element by its index. \par \textbf{Where to use:} Use for Dijkstra's, Prim's, A* search, and any algorithm requiring decrease-key or increase-key operations on specific elements.

\begin{lstlisting}
dsa::indexed_min_pq<int> pq(10);
pq.push(0, 5);
pq.push(1, 3);
pq.push(2, 7);
pq.decrease_key(2, 1);
int idx = pq.pop(); // 2 (min key = 1)
\end{lstlisting}

\subsubsection{Min-Max Heap}
\textbf{Header:} code/include/min\_max\_heap.h\\
\textbf{Class:} dsa::min\_max\_heap$<$T$>$\\

Double-ended priority queue. Supports efficient access to both minimum and maximum elements simultaneously.

\textbf{Why:} You need access to BOTH minimum and maximum simultaneously. \par \textbf{Where to use:} Use for sliding-window min/max and other double-ended priority-queue workloads. Choose based on measured behavior when comparing it with two separate heaps.

\begin{lstlisting}
dsa::min_max_heap<int> mmh;
mmh.push(3); mmh.push(1); mmh.push(5);
mmh.min(); // 1
mmh.max(); // 5
mmh.pop_min(); // removes 1
\end{lstlisting}

\subsubsection{Interval Heap}
\textbf{Header:} code/include/interval\_heap.h\\
\textbf{Class:} dsa::interval\_heap$<$T$>$\\

Double-ended priority queue using two sorted sets, storing interval pairs (lo, hi).

\textbf{Why:} You need a double-ended priority queue that stores interval pairs [lo, hi]. \par \textbf{Where to use:} Use for interval scheduling, deadline management, and event processing where you track both start and end times.

\begin{lstlisting}
dsa::interval_heap<int> ih;
ih.push(2, 8); // interval [2,8]
ih.push(1, 5); // interval [1,5]
ih.min(); // 1
ih.max(); // 8
\end{lstlisting}

\subsubsection{Binomial Heap}
\textbf{Header:} code/include/binomial\_heap.h\\
\textbf{Class:} dsa::BinomialHeap$<$T$>$\\

Mergeable heap using binomial trees. Efficient meld (merge) operation unlike binary heaps.

\textbf{Why:} You need amortized merge of two priority queues, which binary heaps can't do efficiently. \par \textbf{Where to use:} Use for k-way merge in external sorting, Dijkstra's/Prim's on dense graphs, and when you frequently merge two priority queues.

\begin{lstlisting}
dsa::BinomialHeap<int> h1, h2;
h1.push(3); h1.push(1);
h2.push(2);
h1.merge(h2);
int top = h1.pop(); // 1
\end{lstlisting}

\subsubsection{Fibonacci Heap}
\textbf{Header:} code/include/fibonacci\_heap.h\\
\textbf{Class:} dsa::FibonacciHeap$<$T$>$\\

Theoretically optimal decrease-key performance. Amortized insert and decrease-key, delete-min.

\textbf{Why:} You need theoretically optimal amortized decrease-key for dense graph algorithms. \par \textbf{Where to use:} Use when theoretical optimality matters: Dijkstra's, Prim's, and computational geometry. Rarely used in practice due to high constant factors.

\begin{lstlisting}
dsa::FibonacciHeap<int> fh;
fh.push(5); fh.push(3); fh.push(7);
int top = fh.pop(); // 3
fh.push(1);
top = fh.pop(); // 1
\end{lstlisting}

\subsubsection{Pairing Heap}
\textbf{Header:} code/include/pairing\_heap.h\\
\textbf{Class:} dsa::PairingHeap$<$T$>$\\

Simple mergeable heap with two-pass meld. Simpler than Fibonacci heaps with comparable practical performance.

\textbf{Why:} You want Fibonacci-heap-like performance with simpler implementation. \par \textbf{Where to use:} Use as a practical alternative to Fibonacci heaps for Dijkstra's, event simulations, and whenever you need merge + decrease-key without the complexity.

\begin{lstlisting}
dsa::PairingHeap<int> ph;
ph.push(3); ph.push(1); ph.push(4);
int top = ph.pop(); // 1
\end{lstlisting}

\subsubsection{Skew Heap}
\textbf{Header:} code/include/skew\_heap.h\\
\textbf{Class:} dsa::skew\_heap$<$T$>$\\

Self-adjusting mergeable heap. No balance information stored; structure adapts through rotations on every operation.

\textbf{Why:} You need a mergeable heap with self-adjusting structure and no balance overhead. \par \textbf{Where to use:} Use for mergeable priority queues when simplicity matters more than worst-case guarantees. Good for shortest path algorithms.

\begin{lstlisting}
dsa::skew_heap<int> sh1, sh2;
sh1.push(3); sh1.push(1);
sh2.push(2);
sh1.meld(std::move(sh2));
int top = sh1.pop(); // 1
\end{lstlisting}

\subsubsection{Leftist Tree}
\textbf{Header:} code/include/leftist\_tree.h\\
\textbf{Class:} dsa::leftist\_tree$<$T$>$\\

Mergeable min-heap with the leftist property: the null-path length of the left child is always $\geq$ the right child's.

\textbf{Why:} You need merge with the leftist property guaranteeing shortest right path. \par \textbf{Where to use:} Use for k-way merge in external sorting, tournament trees, and when you need guaranteed merge with simpler code than binomial/Fibonacci heaps.

\begin{lstlisting}
dsa::leftist_tree<int> lt1, lt2;
lt1.push(3); lt1.push(1);
lt2.push(2);
lt1.meld(std::move(lt2));
auto sorted = lt1.drain_sorted(); // [1, 2, 3]
\end{lstlisting}

% ---- BALANCED TREES ----
\subsection{Balanced Search Trees}

\subsubsection{AVL Tree}
\textbf{Header:} code/include/avl\_tree.h\\
\textbf{Class:} dsa::avl\_tree$<$K,V$>$\\

Height-balanced BST. Balance factor for every node is in \{-1, 0, 1\}. Strictly balanced, giving optimal lookup time.

\textbf{Why:} You need strictly balanced BST with optimal lookup time and can tolerate more rotations on insert/erase. \par \textbf{Where to use:} Use when lookups dominate over insertions. Stricter balance than red-black trees gives faster search but more rotation overhead on updates.

\begin{lstlisting}
dsa::avl_tree<int,std::string> t;
t.insert(5, "five");
t.insert(3, "three");
t.find(5); // "five"
t.inorder(); // [(3,"three"),(5,"five")]
\end{lstlisting}

\subsubsection{Red-Black Tree}
\textbf{Header:} code/include/red\_black\_tree.h\\
\textbf{Class:} dsa::RedBlackTree$<$Key$>$\\

2-3-4 tree isomorphic. Less rigidly balanced than AVL but fewer rotations on insert/erase. This is the basis for std::map and std::set.

\textbf{Why:} You need a general-purpose ordered map with all operations and fewer rotations than AVL. \par \textbf{Where to use:} Use as the default ordered map/set. Fewer rotations than AVL on insert/erase make it better for write-heavy workloads. This is what std::map uses.

\begin{lstlisting}
dsa::RedBlackTree<int> rbt;
rbt.insert(5); rbt.insert(3); rbt.insert(7);
rbt.contains(5); // true
rbt.erase(3);
auto v = rbt.inorder(); // [5, 7]
\end{lstlisting}

\subsubsection{Splay Tree}
\textbf{Header:} code/include/splay\_tree.h\\
\textbf{Class:} dsa::SplayTree$<$Key$>$\\

Self-adjusting BST. Every access splays the element to the root via rotations, exploiting access locality.

\textbf{Why:} Your workload has access locality (recently accessed items are accessed again soon). \par \textbf{Where to use:} Use when access patterns have temporal locality. Splaying brings frequently accessed elements near root, giving amortized performance with good cache behavior.

\begin{lstlisting}
dsa::SplayTree<int> st;
st.insert(5); st.insert(3); st.insert(7);
st.contains(5); // true, splays 5 to root
st.erase(3);
\end{lstlisting}

\subsubsection{Treap}
\textbf{Header:} code/include/treap.h\\
\textbf{Class:} dsa::treap$<$K,V$>$\\

BST ordered by key, heap-ordered by random priority. The randomization eliminates adversarial input. Supports split() and merge() for set operations.

\textbf{Why:} You need a balanced BST with split/merge operations for set operations. \par \textbf{Where to use:} Use for implicit treaps (sequence manipulation), order-statistics, split/merge set operations, and when you want expected performance without worst-case balance.

\begin{lstlisting}
dsa::treap<int> t;
t.insert(5); t.insert(3); t.insert(7);
t.contains(5); // true
auto v = t.inorder(); // [3, 5, 7]
t.erase(3);
\end{lstlisting}

\subsubsection{2-3 Tree}
\textbf{Header:} code/include/two\_three\_tree.h\\
\textbf{Class:} dsa::TwoThreeTree$<$Key$>$\\

All leaves at same depth. Nodes hold 1-2 keys and have 2-3 children. Conceptually simpler than red-black trees.

\textbf{Why:} You want the simplest B-tree variant for learning or for small datasets where disk I/O isn't the bottleneck. \par \textbf{Where to use:} Use as an educational tool for understanding B-trees and disk-based indexing. Good for small in-memory datasets.

\begin{lstlisting}
dsa::TwoThreeTree<int> t23;
t23.insert(5); t23.insert(3); t23.insert(7);
t23.contains(5); // true
t23.erase(3);
auto v = t23.inorder(); // [5, 7]
\end{lstlisting}

\subsubsection{2-3-4 Tree}
\textbf{Header:} code/include/two\_three\_four\_tree.h\\
\textbf{Class:} dsa::two\_three\_four\_tree$<$K,V$>$\\

B-tree of order 4. Isomorphic to red-black trees. Every internal node has 2-4 children.

\textbf{Why:} You want a B-tree of order 4 with simpler insertion than red-black trees. \par \textbf{Where to use:} Use for database indexing and file systems. Isomorphic to red-black trees but with simpler top-down insertion logic.

\begin{lstlisting}
dsa::two_three_four_tree<int, std::string> t234;
t234.insert(5, "five"); t234.insert(3, "three");
auto* val = t234.find(5); // "five"
t234.erase(3);
\end{lstlisting}

\subsubsection{B-Tree}
\textbf{Header:} code/include/b\_tree.h\\
\textbf{Class:} dsa::b\_tree$<$Key,Compare$>$\\

Generalized balanced search tree with configurable order. Its array-backed nodes illustrate the structure, but the implementation is in-memory and does not perform disk I/O.

\textbf{Why:} You want to study a B-tree with configurable node order. \par \textbf{Where to use:} Use as an in-memory teaching implementation; adapting it for disk-based indexing requires page serialization and an external storage layer.

\begin{lstlisting}
dsa::b_tree<int> bt(5); // order 5
bt.insert(5); bt.insert(3); bt.insert(7);
bt.contains(5); // true
auto v = bt.inorder(); // [3, 5, 7]
\end{lstlisting}

\subsubsection{B+ Tree}
\textbf{Header:} code/include/bplus\_tree.h\\
\textbf{Class:} dsa::bplus\_tree$<$K,V,ORDER$>$\\

All data in leaves; leaves linked for efficient range scans. The workhorse of database indexing.

\textbf{Why:} You need both point lookups AND efficient range queries on sorted data. \par \textbf{Where to use:} Use to study the structure used by many database indexes. This implementation is in-memory; production database use also requires persistence, page management, and concurrency control.

\begin{lstlisting}
dsa::bplus_tree<int, std::string> bpt;
bpt.insert(5, "five"); bpt.insert(3, "three");
auto* val = bpt.find(5); // "five"
auto range = bpt.range(1, 6); // [(3,"three"),(5,"five")]
\end{lstlisting}

\subsubsection{(a,b)-Tree}
\textbf{Header:} code/include/ab\_tree.h\\
\textbf{Class:} dsa::ab\_tree$<$K,V,A,B$>$\\

B-tree generalization with configurable min/max children bounds $a$ and $b$. Provides fine-grained control over node sizes.

\textbf{Why:} You need fine-grained control over node sizes with configurable min/max children bounds. \par \textbf{Where to use:} Use when you want to tune the B-tree family parameters for specific page sizes or memory layouts.

\begin{lstlisting}
dsa::ab_tree<int, std::string, 2, 4> abt;
abt.insert(5, "five"); abt.insert(3, "three");
auto* val = abt.find(5); // "five"
\end{lstlisting}

\subsubsection{Scapegoat Tree}
\textbf{Header:} code/include/scapegoat\_tree.h\\
\textbf{Class:} dsa::scapegoat\_tree$<$K$>$\\

Weight-balanced BST. Rebuilds an entire subtree when balance exceeds $\alpha = 0.75$. No balance info stored per node.

\textbf{Why:} You need a self-balancing BST with zero per-node balance metadata. \par \textbf{Where to use:} Use when memory is tight and you can't afford per-node balance fields. Rebuilds subtrees instead of tracking balance.

\begin{lstlisting}
dsa::scapegoat_tree<int, std::string> sgt;
sgt.insert(5, "five"); sgt.insert(3, "three");
auto* val = sgt.find(5); // "five"
sgt.erase(3);
\end{lstlisting}

\subsubsection{AA Tree}
\textbf{Header:} code/include/aa\_tree.h\\
\textbf{Class:} dsa::aa\_tree$<$K,V$>$\\

Simplified red-black tree variant where all red links are right links. Only skew and split operations needed.

\textbf{Why:} You want the simplest self-balancing BST with only skew and split operations. \par \textbf{Where to use:} Use when you want a simpler red-black tree implementation. Only two operations (skew + split) make it easier to teach and implement correctly.

\begin{lstlisting}
dsa::aa_tree<int, std::string> aat;
aat.insert(5, "five"); aat.insert(3, "three");
auto* val = aat.find(5); // "five"
aat.erase(3);
\end{lstlisting}

\subsubsection{Left-Leaning Red-Black Tree}
\textbf{Header:} code/include/llrb\_tree.h\\
\textbf{Class:} dsa::llrb\_tree$<$K,V$>$\\

Sedgewick's 2008 variant. Enforces that all red links lean left, reducing the number of cases in rotations.

\textbf{Why:} You want a red-black tree variant with fewer rotation cases by enforcing left-leaning red links. \par \textbf{Where to use:} Use when you want a cleaner red-black tree implementation. Sedgewick's variant reduces rotation cases from 7 to 2.

\begin{lstlisting}
dsa::llrb_tree<int, std::string> llrb;
llrb.insert(5, "five"); llrb.insert(3, "three");
auto* val = llrb.find(5); // "five"
llrb.erase(3);
\end{lstlisting}

% ---- MULTI-WAY TREES ----
\subsubsection{Radix/Patricia Trie}
\textbf{Header:} code/include/radix\_tree.h\\
\textbf{Class:} dsa::RadixTree$<$V$>$\\

Compressed trie with merged single-child edges. Edges carry string labels instead of single characters.

\textbf{Why:} Your keys share common prefixes and you need efficient prefix-based operations. \par \textbf{Where to use:} Use for IP routing (longest prefix match), autocomplete, spell-checking, and any prefix-based dictionary. Compressed trie merges single-child edges.

\begin{lstlisting}
dsa::RadixTree<int> rt;
rt.insert("hello", 1); rt.insert("help", 2);
auto* val = rt.find("hello"); // 1
auto matches = rt.prefix_search("hel"); // [("hello",1),("help",2)]
\end{lstlisting}

\subsubsection{Suffix Tree}
\textbf{Header:} code/include/suffix\_tree.h\\
\textbf{Class:} dsa::SuffixTree\\

Naive construction for all suffixes of a string. Pattern matching is linear in the text of the matched subtree plus the number of reported occurrences after preprocessing; construction is \(O(n^2)\).

\textbf{Why:} You need to answer multiple substring queries on the same text after preprocessing. \par \textbf{Where to use:} Use for longest common substring, repeated substring detection, and pattern matching when you'll query the same text many times.

\begin{lstlisting}
dsa::SuffixTree st("bananas");
auto pos = st.search("ana"); // {1, 3}
auto lrs = st.longest_repeated_substring(); // "ana"
\end{lstlisting}

\subsubsection{Suffix Array}
\textbf{Header:} code/include/suffix\_array.h\\
\textbf{Class:} dsa::SuffixArray\\

Space-efficient alternative to suffix trees: the sorted array of all suffixes plus the LCP (longest common prefix) array.

\textbf{Why:} You need space-efficient full-text indexing with pattern search capability. \par \textbf{Where to use:} Use when memory is a concern: suffix arrays generally use less space than suffix trees. Good for full-text search and bioinformatics.

\begin{lstlisting}
dsa::SuffixArray sa("banana");
auto pos = sa.search("ana"); // {1, 3}
auto count = sa.count_occurrences("ana"); // 2
\end{lstlisting}

\subsubsection{Interval Tree}
\textbf{Header:} code/include/interval\_tree.h\\
\textbf{Class:} dsa::interval\_tree$<$T$>$\\

Augmented BST for interval stabbing queries. Each node stores the maximum endpoint in its subtree.

\textbf{Why:} You need to find all intervals containing a given point (stabbing query). \par \textbf{Where to use:} Use for scheduling (finding overlapping meetings), network monitoring, and any interval stabbing query. Augmented BST with max endpoint.

\begin{lstlisting}
dsa::interval_tree<int> it;
it.insert(dsa::interval_tree<int>::interval{1, 5, 0});
it.insert(dsa::interval_tree<int>::interval{3, 8, 1});
auto hits = it.query(4); // both intervals contain 4
\end{lstlisting}

% ---- SPATIAL ----
\subsection{Spatial \& Range Structures}

\subsubsection{k-d Tree}
\textbf{Header:} code/include/kd\_tree.h\\
\textbf{Class:} dsa::kd\_tree\\

Multi-dimensional nearest-neighbor and range queries. Partitions space by alternating splitting dimensions.

\textbf{Why:} You need nearest-neighbor or range queries in multi-dimensional space. \par \textbf{Where to use:} Use for k-NN in ML, GIS queries, collision detection in games, and particle simulation. Good for low-to-medium dimensions.

\begin{lstlisting}
dsa::kd_tree<2, double> kd({{0,0,0},{1,1,1},{5,5,2}});
auto [pt, dist] = kd.nearest({1, 1});
auto near = kd.range_query({1, 1}, 2.0);
\end{lstlisting}

\subsubsection{Quadtree}
\textbf{Header:} code/include/quadtree.h\\
\textbf{Class:} dsa::quadtree\\

Recursive 2D spatial partitioning with NW/NE/SW/SE children.

\textbf{Why:} Your data is 2D and you need recursive spatial partitioning. \par \textbf{Where to use:} Use for 2D collision detection, image compression, and geographic indexing. Good when points are uniformly distributed.

\begin{lstlisting}
dsa::quadtree qt({0, 0, 100, 100});
qt.insert(10, 20, 0);
qt.insert(50, 60, 1);
auto near = qt.range_query(10, 20, 5.0);
\end{lstlisting}

\subsubsection{R-Tree}
\textbf{Header:} code/include/r\_tree.h\\
\textbf{Class:} dsa::r\_tree\\

Spatial index for multi-dimensional bounding-box queries. Groups nearby objects using minimum bounding rectangles.

\textbf{Why:} You need to index multi-dimensional objects (not just points) with bounding-box queries. \par \textbf{Where to use:} Use for spatial databases, map applications, and game spatial indexing. Groups nearby objects using minimum bounding rectangles.

\begin{lstlisting}
dsa::r_tree<2> rt;
rt.insert({{1, 2, 0}});
rt.insert({{3, 4, 1}});
auto pts = rt.range_query({{0,0},{5,5}});
\end{lstlisting}

\subsubsection{Orthogonal Range Tree}
\textbf{Header:} code/include/orthogonal\_range\_tree.h\\
\textbf{Class:} dsa::orthogonal\_range\_tree\\

Segment tree on x-axis, sorted-by-y at each node. Efficient for axis-aligned rectangular range queries.

\textbf{Why:} You need to report all points in an axis-aligned rectangular range. \par \textbf{Where to use:} Use for windowing queries in computational geometry, VLSI checking, and 2D range reporting. Segment tree on x, sorted by y.

\begin{lstlisting}
dsa::orthogonal_range_tree<int> ort({{1,2,0},{3,4,1},{5,6,2}});
auto pts = ort.query(1, 5, 2, 6); // rectangle query
\end{lstlisting}

\subsubsection{Segment Tree (Lazy)}
\textbf{Header:} code/include/segment\_tree\_lazy.h\\
\textbf{Class:} dsa::segment\_tree\_lazy\\

Segment tree with lazy propagation for range updates + range queries. Supports sum, min, or max with range add/set.

\textbf{Why:} You need BOTH range updates (add/set) AND range queries on an array. \par \textbf{Where to use:} Use when you need range update AND range query. Lazy propagation defers updates. Essential for competitive programming.

\begin{lstlisting}
dsa::segment_tree_lazy stl({1, 2, 3, 4, 5});
stl.update(1, 3, 10); // add 10 to indices [1,3]
int s = stl.query(0, 4); // sum of all elements
\end{lstlisting}

\subsubsection{Fenwick Tree (BIT)}
\textbf{Header:} code/include/fenwick\_segment.h\\
\textbf{Classes:} dsa::fenwick\_tree, dsa::segment\_tree, dsa::union\_find\\

Binary Indexed Tree for prefix sums. Simpler to implement than segment trees with similar performance.

\textbf{Why:} You need prefix sums with point updates, and want simpler code than segment trees. \par \textbf{Where to use:} Use for counting inversions, prefix sum queries, and 2D range sums. Simpler than segment trees with similar performance.

\begin{lstlisting}
dsa::fenwick_tree<int> ft(10);
ft.add(0, 5); ft.add(1, 3);
int p = ft.prefix_sum(1); // 8
int r = ft.range_sum(0, 1); // 8
\end{lstlisting}

\subsubsection{Union-Find}
\textbf{Header:} code/include/union\_find.h\\
\textbf{Class:} dsa::union\_find\\

Disjoint-set with path compression + union by size. The inverse Ackermann function $\alpha(n)$ is effectively constant.

\textbf{Why:} You need to efficiently merge sets and check if two elements are in the same set. \par \textbf{Where to use:} Use for Kruskal's MST, connected components in dynamic graphs, cycle detection, and percolation. Amortized time is \(O(\alpha(n))\) per operation.

\begin{lstlisting}
dsa::union_find uf(5);
uf.unite(0, 1); // merge sets
uf.unite(1, 2);
uf.connected(0, 2); // true
uf.count(); // 3 components
\end{lstlisting}

\subsubsection{van Emde Boas Tree}
\textbf{Header:} code/include/veb\_tree.h\\
\textbf{Class:} dsa::veb\_tree\\

Near-constant time for integer operations when the universe U is bounded.

\textbf{Why:} You need predecessor/successor on bounded integer keys. \par \textbf{Where to use:} Use when keys are bounded integers and you need \(O(\log\log U)\) operations. Space grows with the universe size \(U\).

\begin{lstlisting}
dsa::veb_tree<256> veb;
veb.insert(5); veb.insert(3); veb.insert(7);
veb.successor(4); // 5
veb.predecessor(6); // 5
veb.erase(3);
\end{lstlisting}

\subsubsection{X-Fast Trie}
\textbf{Header:} code/include/x\_fast\_trie.h\\
\textbf{Class:} dsa::x\_fast\_trie\\

Trie-based structure achieving van Emde Boas-like performance with less space.

\textbf{Why:} You want vEB-like performance with less space. \par \textbf{Where to use:} Use as a space-efficient alternative to van Emde Boas trees for integer predecessor/successor queries.

\begin{lstlisting}
dsa::x_fast_trie<256> xft;
xft.insert(5); xft.insert(3); xft.insert(7);
xft.successor(4); // 5
xft.predecessor(6); // 5
xft.erase(3);
\end{lstlisting}

\subsubsection{Merkle Tree}
\textbf{Header:} code/include/merkle\_tree.h\\
\textbf{Class:} dsa::merkle\_tree$<$T$>$\\

Hash tree for data integrity verification. Each leaf hashes a data block; each internal node hashes its children.

\textbf{Why:} You need to verify data integrity efficiently across distributed systems. \par \textbf{Where to use:} Use for blockchain, distributed databases, Git, and any system requiring tamper-evident data verification. A proof for one leaf is logarithmic in the number of leaves.

\subsubsection{Bloom Filter}
\textbf{Header:} code/include/bloom\_filter.h\\
\textbf{Class:} dsa::bloom\_filter$<$Bits,K$>$\\

Probabilistic set membership. Space-efficient with tunable false-positive rate. No false negatives.

\textbf{Why:} You need fast set membership testing with low memory, and can tolerate false positives. \par \textbf{Where to use:} Use when you need to filter out definitely-not-present items quickly: URL blacklists, database bloom filters, cache filtering. Never has false negatives.

\begin{lstlisting}
dsa::bloom_filter<> bf;
bf.insert("hello");
bf.might_contain("hello"); // true
bf.might_contain("world"); // maybe false
\end{lstlisting}

%% ============================================================
%% PART III: ALGORITHM DESIGN
%% ============================================================

\section{Part III --- Algorithm Design Paradigms}

% ---- SORTING ----
\subsection{Sorting}
\textbf{Header:} code/include/sorting.h

12 sorting algorithms covering comparison-based, non-comparison, and specialized approaches.

\subsubsection{Insertion Sort}
\textbf{Function:} dsa::insertion\_sort(span$<$T$>$ a)\\

Builds the sorted array one element at a time by inserting each element into its correct position among the already-sorted prefix.

\textbf{Why:} Insertion sort is acceptable for small n, or data is nearly sorted (best case). \par \textbf{Where to use:} Use for small arrays (n < 20), nearly-sorted data, as the base case in hybrid sorts, and online sorting. Best when data is almost sorted.

\begin{lstlisting}
std::vector<int> a = {5, 3, 1, 4, 2};
dsa::insertion_sort(a);
// a = {1, 2, 3, 4, 5}
\end{lstlisting}

\subsubsection{Merge Sort}
\textbf{Function:} dsa::merge\_sort(span$<$T$>$ a)\\

Divide-and-conquer: split in half, recursively sort each half, then merge. Stable sort with guaranteed \(O(n\log n)\) time.

\textbf{Why:} You need guaranteed and/or stable sorting and can afford extra space. \par \textbf{Where to use:} Use when you need stability (preserving relative order of equal elements), guaranteed \(O(n\log n)\) time, or external sorting (data larger than memory).

\begin{lstlisting}
std::vector<int> a = {5, 3, 1, 4, 2};
dsa::merge_sort(a);
// a = {1, 2, 3, 4, 5}
\end{lstlisting}

\subsubsection{Quick Sort}
\textbf{Function:} dsa::quick\_sort(span$<$T$>$ a)\\

Partition around a pivot, recursively sort partitions. Fastest in practice despite worse-case.

\textbf{Why:} You want the fastest average-case sort and can accept worst-case performance (mitigated by randomization). \par \textbf{Where to use:} Use as the default in-memory sort. Fastest in practice due to cache-friendly partitioning. Randomized pivot makes worst-case extremely unlikely.

\begin{lstlisting}
std::vector<int> a = {5, 3, 1, 4, 2};
dsa::quick_sort(a);
// a = {1, 2, 3, 4, 5}
\end{lstlisting}

\subsubsection{3-Way Quick Sort}
\textbf{Function:} dsa::quick\_sort\_3way(span$<$T$>$ a)\\

Dutch National Flag partition: three-way split into less-than, equal-to, and greater-than pivot.

\textbf{Why:} Your data has many duplicate keys and standard quick sort can degrade toward quadratic behavior. \par \textbf{Where to use:} Use when data has many duplicates: database sorting, histograms, and categorical data. Three-way partition avoids repeatedly partitioning equal keys.

\begin{lstlisting}
std::vector<int> a = {5, 3, 3, 4, 3};
dsa::quick_sort_3way(a);
// a = {3, 3, 3, 4, 5}
\end{lstlisting}

\subsubsection{Quickselect}
\textbf{Function:} dsa::quick\_select(span$<$T$>$ a, size\_t k)\\

Finds the k-th smallest element without fully sorting.

\textbf{Why:} You need the k-th smallest element without fully sorting the array. \par \textbf{Where to use:} Use for finding medians, top-k queries, percentiles, and order statistics. Average-case time is \(O(n)\), and the function partitions the input in place.

\begin{lstlisting}
std::vector<int> a = {5, 3, 1, 4, 2};
int median = dsa::quick_select(a, 2); // 3rd smallest = 3
\end{lstlisting}

\subsubsection{Shell Sort}
\textbf{Function:} dsa::shell\_sort(span$<$T$>$ a)\\

Generalization of insertion sort using decreasing gap sequences. Elements far apart are sorted first, then reduced.

\textbf{Why:} You need a simple in-place sort with better practical performance than insertion sort and minimal code. \par \textbf{Where to use:} Use in embedded systems, medium arrays (n ~ 1000), and when code simplicity matters more than optimal performance.

\begin{lstlisting}
std::vector<int> a = {5, 3, 1, 4, 2};
dsa::shell_sort(a);
// a = {1, 2, 3, 4, 5}
\end{lstlisting}

\subsubsection{Bubble Sort}
\textbf{Function:} dsa::bubble\_sort(span$<$T$>$ a)\\

Repeatedly swap adjacent out-of-order elements. Stops early if no swaps occur.

\textbf{Why:} You need the simplest possible sort, or want to detect if data is already sorted. \par \textbf{Where to use:} Use for teaching, detecting sorted arrays (\(O(n)\) best case), and tiny datasets. Never use in production for large data.

\begin{lstlisting}
std::vector<int> a = {5, 3, 1, 4, 2};
dsa::bubble_sort(a);
// a = {1, 2, 3, 4, 5}
\end{lstlisting}

\subsubsection{Bottom-Up Merge Sort}
\textbf{Function:} dsa::merge\_sort\_bottom\_up(span$<$T$>$ a)\\

Iterative merge sort avoiding recursion overhead. Merges subarrays of increasing size.

\textbf{Why:} You need merge sort without recursion overhead. \par \textbf{Where to use:} Use in embedded systems where recursion is costly, and for linked list sorting without random access.

\begin{lstlisting}
std::vector<int> a = {5, 3, 1, 4, 2};
dsa::merge_sort_bottom_up(a);
// a = {1, 2, 3, 4, 5}
\end{lstlisting}

\subsubsection{Counting Sort}
\textbf{Function:} dsa::counting\_sort(span$<$int$>$ a)\\

Non-comparison sort using an auxiliary count array.

\textbf{Why:} Your data consists of integers with a known, small range. \par \textbf{Where to use:} Use when the key range \(k\) is small relative to the input size. Its \(O(n+k)\) time can beat comparison sorts. Also used as a subroutine in radix sort.

\begin{lstlisting}
std::vector<int> a = {4, 2, 5, 2, 1};
dsa::counting_sort(a);
// a = {1, 2, 2, 4, 5}
\end{lstlisting}

\subsubsection{Radix Sort (LSD)}
\textbf{Function:} dsa::radix\_sort(span$<$int$>$ a)\\

Least-significant-digit first. Sorts by each digit position using counting sort.

\textbf{Why:} You're sorting fixed-length integers and want performance. \par \textbf{Where to use:} Use for fixed-length keys: phone numbers, SSNs, dates, where d = number of digits. Faster than comparison sorts for bounded data.

\begin{lstlisting}
std::vector<int> a = {170, 45, 75, 90, 802};
dsa::radix_sort(a);
// a = {45, 75, 90, 170, 802}
\end{lstlisting}

\subsubsection{Radix Sort (MSD for Strings)}
\textbf{Function:} dsa::radix\_sort\_msd(span$<$string$>$ a)\\

Most-significant-digit first, recursive. Naturally handles variable-length strings.

\textbf{Why:} You're sorting variable-length strings by their lexicographic order. \par \textbf{Where to use:} Use for sorting strings: dictionaries, URLs, file paths. MSD radix sort naturally handles variable-length strings.

\begin{lstlisting}
std::vector<std::string> a = {"banana", "apple", "cherry"};
dsa::radix_sort_msd(a);
// a = {"apple", "banana", "cherry"}
\end{lstlisting}

\subsubsection{Bucket Sort}
\textbf{Function:} dsa::bucket\_sort(span$<$float$>$ a)\\

Distributes elements into buckets, sorts each bucket, concatenates.

\textbf{Why:} Your data is uniformly distributed floating-point numbers in [0,1). \par \textbf{Where to use:} Use when data is uniformly distributed: floating-point numbers, histogram equalization. Expected time is \(O(n)\) for a suitable distribution.

\begin{lstlisting}
std::vector<float> a = {0.7f, 0.1f, 0.4f, 0.2f};
dsa::bucket_sort(a);
// a = {0.1f, 0.2f, 0.4f, 0.7f}
\end{lstlisting}

% ---- GREEDY ----
\subsection{Greedy Algorithms}
\textbf{Header:} code/include/greedy.h

\subsubsection{Fractional Knapsack}
\textbf{Function:} dsa::fractional\_knapsack(capacity, items)\\

Greedy by value-to-weight ratio. Takes fractions of items to maximize value within capacity.

\textbf{Why:} You can take fractions of items to maximize value (items are divisible). \par \textbf{Where to use:} Use when items are divisible: bandwidth sharing, CPU time slicing, commodity trading. Greedy by value/weight ratio gives optimal solution.

\begin{lstlisting}
std::vector<dsa::frac_item> items =
 {{10,60},{20,100},{30,120}};
double val = dsa::fractional_knapsack(50, items);
// val = 240.0
\end{lstlisting}

\subsubsection{Huffman Coding}
\textbf{Function:} dsa::huffman\_encode(freq\_table)\\

Builds an optimal prefix-free binary code tree. Frequent characters get shorter codes.

\textbf{Why:} You need optimal prefix-free encoding for known symbol frequencies. \par \textbf{Where to use:} Use for entropy encoding in compression: ZIP, JPEG, H.264. Gives optimal prefix-free codes. Often combined with LZ77/LZ78.

\begin{lstlisting}
std::vector<std::pair<char,int>> freq =
 {{'a',5},{'b',9},{'c',12}};
auto codes = dsa::huffman_encode(freq);
// codes: [('a',"00"), ('b',"01"), ('c',"1")]
\end{lstlisting}

% ---- DIVIDE AND CONQUER ----
\subsection{Divide \& Conquer}
\textbf{Header:} code/include/divide\_conquer.h

\subsubsection{Closest Pair of Points}
\textbf{Function:} dsa::closest\_pair(points)\\

Divides points by x-coordinate, recursively finds closest pair in each half, then checks the strip around the dividing line.

\textbf{Why:} You need to find the closest pair among \(n\) points in 2D. \par \textbf{Where to use:} Use for facility location, collision detection, and clustering. The divide-and-conquer implementation runs in \(O(n\log n)\) time.

\begin{lstlisting}
std::vector<std::pair<double,double>> pts =
 {{0,0},{1,1},{5,5},{6,6}};
auto [d, p1, p2] = dsa::closest_pair(pts);
// d = 1.414.. (distance between (0,0) and (1,1))
\end{lstlisting}

\subsubsection{Strassen's Matrix Multiply}
\textbf{Function:} dsa::strassen\_multiply(A, B)\\

Reduces 8 recursive multiplications to 7 at the cost of extra additions. Asymptotically faster than the naive \(O(n^3)\) algorithm.

\textbf{Why:} You need to multiply large square matrices faster asymptotically than the conventional \(O(n^3)\) algorithm. \par \textbf{Where to use:} Use for large matrix multiplication where \(n\) is large enough to offset overhead. Reduces eight recursive multiplications to seven, but this implementation operates on double-precision matrices.

\begin{lstlisting}
std::vector<std::vector<double>> A =
 {{1,2},{3,4}};
std::vector<std::vector<double>> B =
 {{5,6},{7,8}};
auto C = dsa::strassen_multiply(A, B);
// C = {{19,22},{43,50}}
\end{lstlisting}

% ---- DYNAMIC PROGRAMMING ----
\subsection{Dynamic Programming}
\textbf{Header:} code/include/dynamic\_programming.h

\subsubsection{0/1 Knapsack}
\textbf{Function:} dsa::knapsack\_01(capacity, items)\\

Classic DP: fill a 2D table where dp[i][w] = best value using first i items with capacity w. Each item taken at most once.

\textbf{Why:} You have items with weights/values and a capacity constraint, and items can't be split. \par \textbf{Where to use:} Use for budget optimization, cargo loading, and investment selection. Pseudo-polynomial time. Use for small W or exact solutions needed.

\begin{lstlisting}
std::vector<std::pair<int,int>> items =
 {{2,3},{3,4},{4,5},{5,6}};
int val = dsa::knapsack_01(8, items);
// val = 10
\end{lstlisting}

\subsubsection{Longest Common Subsequence}
\textbf{Function:} dsa::lcs(a, b)\\

Finds the longest subsequence present in both strings (not necessarily contiguous).

\textbf{Why:} You need the longest common subsequence (not substring) between two sequences. \par \textbf{Where to use:} Use for diff tools (git diff), version control, DNA alignment, and plagiarism detection. Shows how two sequences differ.

\begin{lstlisting}
auto s = dsa::lcs("ABCBDAB", "BDCAB");
// s = "BCAB" or "BDAB"
\end{lstlisting}

\subsubsection{Edit Distance}
\textbf{Function:} dsa::edit\_distance(a, b)\\

Minimum number of insertions, deletions, and substitutions to transform one string into another.

\textbf{Why:} You need the minimum number of edits to transform one string into another. \par \textbf{Where to use:} Use for spell checking, DNA alignment, speech recognition, and plagiarism detection. Foundation for fuzzy matching.

\begin{lstlisting}
int dist = dsa::edit_distance("kitten", "sitting");
// dist = 3
\end{lstlisting}

\subsubsection{Longest Increasing Subsequence}
\textbf{Function:} dsa::lis\_length(a)\\

Finds the length of the longest strictly increasing subsequence using patience sorting.

\textbf{Why:} You need the length of the longest strictly increasing subsequence. \par \textbf{Where to use:} Use for sequence analysis and as the core of patience sorting. The implementation runs in \(O(n\log n)\) time.

\begin{lstlisting}
std::vector<int> a = {10,9,2,5,3,7,101,18};
auto len = dsa::lis_length(a);
// len = 4 (subsequence: 2,3,7,101)
\end{lstlisting}

% ---- BACKTRACKING ----
\subsection{Backtracking}
\textbf{Header:} code/include/backtracking.h

\subsubsection{N-Queens}
\textbf{Function:} dsa::n\_queens(n)\\

Place N queens on an N x N board so no two attack each other. Uses column and diagonal tracking for efficient pruning.

\textbf{Why:} You need to find all valid configurations satisfying constraints (classic CSP). \par \textbf{Where to use:} Use for constraint satisfaction, puzzle solving, circuit board layout, and as a teaching example for backtracking with pruning.

\begin{lstlisting}
auto sols = dsa::n_queens(4);
// sols.size() = 2
// sols[0] = {".Q.","..Q","Q..",".Q."}
\end{lstlisting}

\subsubsection{Subset Sum}
\textbf{Function:} dsa::subset\_sum(nums, target)\\

Determines if any subset of integers sums to a given target. Uses DP table (not pure backtracking).

\textbf{Why:} You need to determine if any subset sums to a target value. \par \textbf{Where to use:} Use for partition problems, fair division, and cryptographic applications. Pseudo-polynomial time.

\begin{lstlisting}
bool ok = dsa::subset_sum({3, 34, 4, 12, 5, 2}, 9);
// ok = true (4 + 5 = 9)
\end{lstlisting}

\subsubsection{Traveling Salesman}
\textbf{Function:} dsa::tsp(dist)\\

Exact TSP solution via bitmask dynamic programming. Visits all cities exactly once and returns to start.

\textbf{Why:} You need the optimal tour visiting all cities exactly once (small n). \par \textbf{Where to use:} Use for small instances (n <= 20) of routing problems. For larger instances, use heuristics (nearest neighbor, 2-opt, genetic algorithms).

\begin{lstlisting}
std::vector<std::vector<int>> d =
 {{0,10,15,20},{10,0,35,25},
 {15,35,0,30},{20,25,30,0}};
int cost = dsa::tsp(d);
// cost = 80
\end{lstlisting}

% ---- PROBABILISTIC ----
\subsection{Probabilistic Structures}
\textbf{Header:} code/include/probabilistic.h

\subsubsection{Reservoir Sampling}
\textbf{Function:} dsa::reservoir\_sample(k, stream)\\

Selects k items uniformly at random from a stream of unknown length. Each item has equal probability k/n of being selected.

\textbf{Why:} You need k random items from a stream of unknown length in one pass. \par \textbf{Where to use:} Use for streaming random sampling, database sampling, A/B testing, and Monte Carlo. Each item has equal probability k/n.

\begin{lstlisting}
std::vector<int> stream = {1,2,3,4,5,6,7,8,9,10};
auto sample = dsa::reservoir_sample(3, stream);
// sample.size() == 3, uniformly random
\end{lstlisting}

\subsubsection{HyperLogLog}
\textbf{Class:} dsa::hyperloglog\\

Probabilistic cardinality estimator. Uses the observation that hash values with many leading zeros are rare.

\textbf{Why:} You need to count distinct elements in a stream with sublinear space. \par \textbf{Where to use:} Use for unique visitor counting, network flow monitoring, and SELECT DISTINCT estimation. per element, tunable accuracy.

\begin{lstlisting}
dsa::hyperloglog hll(10);
hll.add("hello"); hll.add("world");
hll.add("hello"); // duplicate
hll.count(); // ~2 (approximate cardinality)
\end{lstlisting}

\subsubsection{Count-Min Sketch}
\textbf{Class:} dsa::count\_min\_sketch\\

Probabilistic frequency estimator. Over-estimates but never under-estimates counts.

\textbf{Why:} You need to estimate frequencies in a stream, and can tolerate over-estimation. \par \textbf{Where to use:} Use for network flow counting, stream frequency estimation, and heavy-hitter detection. Never under-estimates, may over-estimate.

\begin{lstlisting}
dsa::count_min_sketch cms(1024, 5);
cms.add("hello", 3);
cms.add("world", 2);
cms.estimate("hello"); // >= 3
\end{lstlisting}

%% ============================================================
%% PART IV: STRINGS & SPECIALIZED
%% ============================================================

\section{Part IV --- Strings \& Specialized Topics}

% ---- EXACT MATCHING ----
\subsection{Exact String Matching}
\textbf{Header:} code/include/exact\_matching.h

Five pattern matching algorithms. All return vector$<$size\_t$>$ of match start positions.

\subsubsection{Morris-Pratt}
\textbf{Function:} dsa::morris\_pratt\_search(text, pattern)\\

Precomputes prefix function to avoid redundant comparisons. The text pointer never moves backward.

\textbf{Why:} You need pattern matching with simple prefix function. \par \textbf{Where to use:} Use as the conceptual stepping stone to KMP. Text pointer never moves backward. Simpler than KMP but skips fewer characters.

\begin{lstlisting}
auto pos = dsa::morris_pratt_search(
 "ABABDABACDABABCABAB", "ABABCABAB");
// pos = {9}
\end{lstlisting}

\subsubsection{Knuth-Morris-Pratt (KMP)}
\textbf{Function:} dsa::kmp\_search(text, pattern)\\

Optimized prefix function uses pattern lookahead to skip more characters on mismatch.

\textbf{Why:} You need the standard pattern matching algorithm. \par \textbf{Where to use:} Use as the default linear-time pattern matcher. KMP's prefix function skips more characters than Morris-Pratt on mismatch.

\begin{lstlisting}
auto pos = dsa::kmp_search(
 "ABABDABACDABABCABAB", "ABABCABAB");
// pos = {9}
\end{lstlisting}

\subsubsection{Karp-Rabin}
\textbf{Function:} dsa::karp\_rabin(pattern, base, mod)\\

Uses rolling hash to compare pattern against text windows in . Only verifies on hash match.

\textbf{Why:} You want average-case with simple rolling hash, and can accept worst case. \par \textbf{Where to use:} Use for plagiarism detection, multiple pattern search, and when simplicity matters. Rolling hash gives window comparison.

\begin{lstlisting}
auto matches = dsa::karp_rabin("ABCABC", "ABC");
// matches = {0, 3}
\end{lstlisting}

\subsubsection{Galil-Seiferas}
\textbf{Function:} dsa::galil\_seiferas\_search(text, pattern)\\

Uses critical factorization for linear-time matching with skip optimization.

\textbf{Why:} You need optimal matching with skip optimization via critical factorization. \par \textbf{Where to use:} Use for long repeated patterns where skip optimization matters. Theoretically optimal but rarely used in practice.

\begin{lstlisting}
auto pos = dsa::galil_seiferas_search("ABCABC", "ABC");
// pos = {0, 3}
\end{lstlisting}

\subsubsection{Boyer-Moore-Horspool}
\textbf{Function:} dsa::boyer\_moore\_horspool(pattern)\\

Simplified Boyer-Moore using only bad-character shift. Compares right-to-left, skips large portions of text.

\textbf{Why:} You want the fastest practical string search for large alphabets. \par \textbf{Where to use:} Use as the default practical string search. Bad-character shift skips large portions of text. Sublinear average-case.

\begin{lstlisting}
dsa::boyer_moore_horspool bm("needle");
auto pos = bm.search("haystack needle here");
// pos = {8}
\end{lstlisting}

% ---- STRING SEARCH ----
\subsection{String Search \& Pattern Matching}
\textbf{Header:} code/include/string\_search.h

\subsubsection{Z-Algorithm}
\textbf{Function:} dsa::z\_function(s)\\

For each position i, computes z[i] = length of longest substring starting at i that matches the prefix.

\textbf{Why:} Pattern matching (concatenate pattern + \$ + text, find z-values equal to pattern length), shortest unique substring, and run-length detection.

\begin{lstlisting}
auto z = dsa::z_function("aabxaab");
auto pos = dsa::z_search("aabxaabxcaab", "aab");
// pos = {0, 9}
\end{lstlisting}

\subsubsection{Boyer-Moore (Good Suffix)}
\textbf{Function:} dsa::boyer\_moore(text, pattern)\\

Uses both bad-character and good-suffix heuristics for maximum skip distance.

\textbf{Why:} You want the fastest practical string search using both bad-character and good-suffix heuristics. \par \textbf{Where to use:} Use as the fastest practical string search. Both heuristics give maximum skip distance. Sublinear average-case.

\begin{lstlisting}
dsa::boyer_moore bm("needle");
auto pos = bm.search("haystack needle here");
// pos = {8}
\end{lstlisting}

\subsubsection{Manacher's Palindrome}
\textbf{Function:} dsa::manacher(s)\\

Finds all odd-length palindromes centered at each position in linear time using symmetry exploitation.

\textbf{Why:} You need to find all palindromic substrings in linear time. \par \textbf{Where to use:} Use for longest palindromic substring, palindrome detection in DNA, and text compression. The linear-time algorithm exploits palindrome symmetry.

\begin{lstlisting}
auto lps = dsa::longest_palindrome("babad");
// lps = "bab" or "aba"
\end{lstlisting}

% ---- APPROXIMATE MATCHING ----
\subsection{Approximate String Matching}
\textbf{Header:} code/include/approximate\_matching.h

\subsubsection{Levenshtein Distance}
\textbf{Function:} dsa::levenshtein(a, b)\\

Minimum edit distance: insertions, deletions, and substitutions.

\textbf{Why:} You need the minimum edit distance (insertions + deletions + substitutions). \par \textbf{Where to use:} Use for spell checking, DNA alignment, speech recognition, and record deduplication. The standard edit distance metric.

\begin{lstlisting}
int dist = dsa::levenshtein("kitten", "sitting");
// dist = 3
\end{lstlisting}

\subsubsection{Hamming Distance}
\textbf{Function:} dsa::hamming\_distance(a, b)\\

Number of positions where corresponding characters differ. Only defined for equal-length strings.

\textbf{Why:} You need to count positions where two equal-length strings differ. \par \textbf{Where to use:} Use for fixed-length string comparison: error-correcting codes, hash comparison, DNA mutation counting. Only for equal-length strings.

\begin{lstlisting}
auto dist = dsa::hamming_distance("karolin", "kathrin");
// dist = 3
\end{lstlisting}

\subsubsection{Myers' Bit-Parallel}
\textbf{Function:} dsa::myers\_bit\_parallel(text, pattern, threshold)\\

Simulates the DP edit distance matrix using bit operations. Extremely fast for small patterns and thresholds.

\textbf{Why:} You need approximate pattern matching with small edit distance threshold. \par \textbf{Where to use:} Use for bioinformatics alignment, spell checking, and approximate grep. Bit-parallel simulation is extremely fast for small patterns.

\begin{lstlisting}
auto pos = dsa::myers_approximate(
 "hello world", "helo", 1);
// pos = {0} (distance 1: 'l' missing)
\end{lstlisting}

\subsubsection{Shift-Or (Approximate)}
\textbf{Function:} dsa::shift\_or\_approximate(text, pattern, threshold)\\

Bit-parallel approximate matching using the Shift-Or framework.

\textbf{Why:} You need approximate pattern matching using bit-parallel DP. \par \textbf{Where to use:} Use for fast approximate matching with Hamming distance. Bit-parallel approach processes multiple positions simultaneously.

\begin{lstlisting}
auto pos = dsa::shift_or_approximate(
 "hello world", "helo", 1);
// pos = {0}
\end{lstlisting}

% ---- MULTIPLE MATCHING ----
\subsection{Multiple Pattern Matching}
\textbf{Header:} code/include/multiple\_matching.h

\subsubsection{Aho-Corasick}
\textbf{Function:} dsa::aho\_corasick\_search(text, patterns)\\

Builds an automaton from all patterns simultaneously. Scans text once, finding all pattern occurrences.

\textbf{Why:} You need to find ALL occurrences of MULTIPLE patterns in a single text scan. \par \textbf{Where to use:} Use for multi-pattern matching: antivirus scanning, intrusion detection, bioinformatics. Builds one automaton from all patterns.

\begin{lstlisting}
dsa::aho_corasick ac({"he", "she", "his", "hers"});
auto matches = ac.search("ahishers");
// matches: (1,"his"), (1,"he"), (4,"she"), (4,"he"), (4,"hers")
\end{lstlisting}

\subsubsection{Commentz-Walter}
\textbf{Function:} dsa::commentz\_walter\_search(text, patterns)\\

Multi-pattern search combining Shift-Or with Aho-Corasick concepts.

\textbf{Why:} You need multi-pattern search with variable-length patterns. \par \textbf{Where to use:} Use as an alternative to Aho-Corasick when pattern lengths vary significantly. Combines Shift-Or with Aho-Corasick concepts.

\begin{lstlisting}
dsa::commentz_walter cw({"he", "she", "his", "hers"});
auto matches = cw.search("ahishers");
// matches: (position, pattern_id) for each match
\end{lstlisting}

% ---- TEXT DECOMPOSITION ----
\subsection{Text Decomposition}
\textbf{Header:} code/include/text\_decomposition.h

\subsubsection{Lyndon Factorization}
\textbf{Function:} dsa::lyndon\_factorization(s)\\

Duval's algorithm decomposes a string into Lyndon words (lexicographically strictly smallest rotations).

\textbf{Why:} You need to decompose a string into Lyndon words. \par \textbf{Where to use:} Use for minimal string rotation (Booth's), suffix array construction (DC3/SA-IS), and string periodicity analysis.

\begin{lstlisting}
auto factors = dsa::lyndon_factorization("banana");
// factors: "b" "an" "an" "a" (Lyndon factorization)
bool is_ld = dsa::is_lyndon_word("an");
// is_ld = true
\end{lstlisting}

\subsubsection{Critical Factorization}
\textbf{Function:} dsa::critical\_factorization(s)\\

Finds the critical index that enables optimal two-way string matching.

\textbf{Why:} You need the critical index for optimal two-way string matching. \par \textbf{Where to use:} Use as preprocessing for Crochemore-Perrin two-way algorithm. Enables matching with left-to-right text scan.

\begin{lstlisting}
auto [pos, period] = dsa::critical_factorization("abcabc");
// pos = 3, period = 3
\end{lstlisting}

% ---- COMPRESSION ----
\subsection{String Compression}
\textbf{Header:} code/include/string\_compression.h

Four classic compression algorithms with encode/decode pairs.

\subsubsection{Run-Length Encoding}
\textbf{Functions:} dsa::rle\_encode(s), dsa::rle\_decode(s)\\

Replaces consecutive repeated characters with character + count.

\textbf{Why:} Your data has long runs of repeated characters. \par \textbf{Where to use:} Use for run-length data: fax, BMP images, sensor data. Simple but only effective when runs are long. Often combined with other methods.

\begin{lstlisting}
auto enc = dsa::rle_encode("AAABBBCCD");
// enc = "A3B3C2D1"
auto dec = dsa::rle_decode(enc);
// dec = "AAABBBCCD"
\end{lstlisting}

\subsubsection{LZ77}
\textbf{Functions:} dsa::lz77\_encode(s), dsa::lz77\_decode(tokens)\\

Sliding-window dictionary compression. References previous occurrences via (offset, length) tuples.

\textbf{Why:} You need sliding-window dictionary compression. \par \textbf{Where to use:} Use as the foundation of DEFLATE (gzip, PNG, ZIP). References previous occurrences via (offset, length) tuples.

\begin{lstlisting}
auto tokens = dsa::lz77_encode("ABCABCABC");
auto orig = dsa::lz77_decode(tokens);
// orig == "ABCABCABC"
\end{lstlisting}

\subsubsection{LZ78}
\textbf{Functions:} dsa::lz78\_encode(s), dsa::lz78\_decode(tokens)\\

Dictionary-based compression building an explicit dictionary during encoding.

\textbf{Why:} You need dictionary-based compression with explicit dictionary building. \par \textbf{Where to use:} Use for LZW variant compression: GIF, Unix compress. Builds dictionary explicitly during encoding.

\begin{lstlisting}
auto tokens = dsa::lz78_encode("TOBEORNOTTOBE");
auto orig = dsa::lz78_decode(tokens);
// orig == "TOBEORNOTTOBE"
\end{lstlisting}

\subsubsection{LZW}
\textbf{Functions:} dsa::lzw\_encode(s), dsa::lzw\_decode(codes)\\

Adaptive dictionary compression. Dictionary is built implicitly during both encoding and decoding.

\textbf{Why:} You need adaptive dictionary compression where dictionary is built during both encode and decode. \par \textbf{Where to use:} Use for GIF, TIFF, PDF compression. Dictionary is built implicitly, no need to transmit it separately.

\begin{lstlisting}
auto codes = dsa::lzw_encode("TOBEORNOTTOBE");
auto orig = dsa::lzw_decode(codes);
// orig == "TOBEORNOTTOBE"
\end{lstlisting}

% ---- BWT ----
\subsection{Burrows-Wheeler Transform}
\textbf{Header:} code/include/bwt.h\\
\textbf{Functions:} dsa::bwt\_encode(s), dsa::bwt\_decode(s)\\

Reorders a string so that similar characters cluster together, enabling high compression ratios.

\textbf{Why:} You need to reorder a string so similar characters cluster together for better compression. \par \textbf{Where to use:} Use as preprocessing for bzip2 compression. Enables high compression ratios by clustering similar characters.

\begin{lstlisting}
auto encoded = dsa::bwt_encode("banana");
auto decoded = dsa::bwt_decode(encoded);
// decoded == "banana"
\end{lstlisting}

\subsubsection{FM-Index}
\textbf{Class:} dsa::fm\_index\\

Compressed full-text index using BWT. Supports substring search and counting without decompression.

\textbf{Why:} You need to search text using its BWT representation. \par \textbf{Where to use:} Use for full-text search in applications such as genome indexing. This implementation stores auxiliary tables and supports counting and locating matches.

\begin{lstlisting}
dsa::fm_index fm("banana");
auto count = fm.count("ana"); // 2
auto pos = fm.search("ana"); // {1, 3}
\end{lstlisting}

% ---- SUFFIX ARRAYS ----
\subsection{Suffix Array Search}
\textbf{Header:} code/include/suffix\_array\_search.h

\subsubsection{Kasai LCP}
\textbf{Function:} dsa::kasai\_lcp(s, sa)\\

Constructs the LCP (longest common prefix) array from a suffix array in linear time.

\textbf{Why:} You need the LCP array from a suffix array. \par \textbf{Where to use:} Use for longest repeated substring, distinct substring counting, and as input to many string algorithms. Kasai's algorithm constructs it in linear time.

\begin{lstlisting}
std::string s = "banana";
std::vector<size_t> sa = {5, 3, 1, 0, 4, 2};
auto lcp = dsa::kasai_lcp(s, sa);
// lcp = {-, 1, 3, 0, 0, 2} (lcp[0] undefined)
\end{lstlisting}

\subsubsection{SA Search}
\textbf{Function:} dsa::sa\_search(sa, s, pattern)\\

Binary search on suffix array to find pattern occurrences.

\textbf{Why:} You need to find pattern occurrences using binary search on a suffix array. \par \textbf{Where to use:} Use for full-text search on suffix arrays. Query time is \(O(m\log n + occ)\), where \(m\) is pattern length and \(occ\) is the number of matches.

\begin{lstlisting}
std::string s = "banana";
std::vector<size_t> sa = {5, 3, 1, 0, 4, 2};
auto pos = dsa::sa_search(sa, s, "ana");
// pos = {1, 3}
\end{lstlisting}

% ---- NFA ----
\subsection{NFA Regex Matching}
\textbf{Header:} code/include/nfa.h\\
\textbf{Function:} dsa::nfa\_match(pattern, text)\\

Thompson's NFA construction for regular expression matching. Supports *, +, ?, |, and ().

\textbf{Why:} You need regular expression matching with *, +, ?, |, and (). \par \textbf{Where to use:} Use for regex matching: input validation, log filtering, web scraping. Thompson's NFA construction gives matching.

\begin{lstlisting}
bool ok = dsa::nfa_match("(a|b)*abb", "aababb");
// ok = true
\end{lstlisting}

% ---- GRAPH ----
\subsection{Graph Algorithms}
\textbf{Header:} code/include/graph.h\\
\textbf{Class:} dsa::graph$<$Weight$>$

Adjacency-list graph with both directed and undirected edge support.

\subsubsection{DFS / BFS}
\textbf{Functions:} g.dfs(start, visitor), g.bfs(start, visitor)\\

DFS uses a stack (recursive); BFS uses a queue. Both visit every reachable vertex exactly once.

\textbf{Applications (DFS):} Topological sorting, cycle detection, connected components, maze solving, path finding in puzzles, and backtracking search.

\textbf{Applications (BFS):} Shortest paths in unweighted graphs, level-order tree traversal, peer-to-peer network discovery, web crawlers, and social network degree-of-separation.

\begin{lstlisting}
dsa::graph<int> g(4);
g.add_edge(0, 1); g.add_edge(1, 2); g.add_edge(2, 3);
g.dfs(0, [](int v) { /* visit v */ });
g.bfs(0, [](int v) { /* visit v */ });
\end{lstlisting}

\subsubsection{Connected Components}
\textbf{Function:} g.connected\_components()\\

Assigns a component ID to each vertex using BFS. Vertices in the same component are mutually reachable.

\textbf{Why:} You need to find all connected components in an undirected graph. \par \textbf{Where to use:} Use for friend group analysis, image processing, network reliability, and dependency analysis. Assigns component IDs via BFS.

\begin{lstlisting}
dsa::graph<int> g(4);
g.add_edge(0, 1); g.add_edge(2, 3);
auto cc = g.connected_components();
// cc = {0,0,1,1} (two components)
\end{lstlisting}

\subsubsection{Dijkstra}
\textbf{Function:} g.dijkstra(start)\\

Finds shortest paths from a source vertex with non-negative edge weights using a min-priority queue.

\textbf{Why:} You need shortest paths from a source with non-negative edge weights. \par \textbf{Where to use:} Use for GPS, IP routing, map applications, and robot path planning. With a binary heap, the usual bound is \(O((V+E)\log V)\).

\begin{lstlisting}
dsa::graph<int> g(5);
g.add_edge(0, 1, 10);
g.add_edge(1, 3, 10);
auto dist = g.dijkstra(0);
// dist[3] = 20
\end{lstlisting}

\subsubsection{Bellman-Ford}
\textbf{Function:} g.bellman\_ford(start)\\

Shortest paths from source, handling negative edge weights. Detects negative cycles.

\textbf{Why:} You need shortest paths and the graph may have negative edge weights. \par \textbf{Where to use:} Use when negative edges exist: arbitrage detection and currency exchange. The standard running time is \(O(VE)\); negative-cycle handling should be checked in the returned result/API contract.

\begin{lstlisting}
dsa::graph<int> g(3);
g.add_edge(0, 1, 4); g.add_edge(1, 2, -2);
auto dist = g.bellman_ford(0);
// dist = {0, 4, 2}
\end{lstlisting}

\subsubsection{Floyd-Warshall}
\textbf{Function:} g.floyd\_warshall()\\

All-pairs shortest paths via dynamic programming. Simple triple-nested loop.

\textbf{Why:} You need all-pairs shortest paths in a dense graph. \par \textbf{Where to use:} Use for small or dense graphs and transitive-closure-style problems. The triple loop takes \(O(V^3)\) time and \(O(V^2)\) space.

\begin{lstlisting}
dsa::graph<int> g(3);
g.add_edge(0, 1, 3); g.add_edge(1, 2, 1);
auto all = g.floyd_warshall();
// all[0][2] = 4
\end{lstlisting}

\subsubsection{Topological Sort}
\textbf{Function:} g.topological\_sort()\\

Linear ordering of vertices in a DAG such that every edge goes from earlier to later.

\textbf{Why:} You need a linear ordering of vertices in a DAG respecting edge constraints. \par \textbf{Where to use:} Use for dependency resolution: build systems, course prerequisites, and task scheduling. The algorithm runs in \(O(V+E)\) time.

\begin{lstlisting}
dsa::graph<int> g(4);
g.add_edge(0, 1); g.add_edge(0, 2); g.add_edge(1, 3);
auto order = g.topological_sort();
// order = {0, 1, 2, 3} or {0, 2, 1, 3}
\end{lstlisting}

\subsubsection{Cycle Detection}
\textbf{Function:} g.has\_cycle()\\

DFS-based detection using three coloring states (Unvisited, InStack, Done).

\textbf{Why:} You need to detect if a directed graph contains a cycle. \par \textbf{Where to use:} Use for deadlock detection, circular dependency checking, and infinite loop detection. DFS-based with three coloring states.

\begin{lstlisting}
dsa::graph<int> g(3);
g.add_edge(0, 1); g.add_edge(1, 2);
g.has_cycle(); // false
g.add_edge(2, 0);
g.has_cycle(); // true
\end{lstlisting}

\subsubsection{Bipartiteness Check}
\textbf{Function:} g.is\_bipartite()\\

2-coloring via BFS. Returns the coloring or empty if the graph is not bipartite.

\textbf{Why:} You need to 2-color a graph or detect odd-length cycles. \par \textbf{Where to use:} Use for matching problems, graph coloring lower bounds, and bipartiteness checking. BFS-based 2-coloring.

\begin{lstlisting}
dsa::graph<int> g(4);
g.add_edge(0, 1); g.add_edge(1, 2); g.add_edge(2, 3);
auto color = g.is_bipartite();
// color = {0, 1, 0, 1} (bipartite)
\end{lstlisting}

\subsubsection{Transitive Closure}
\textbf{Function:} g.transitive\_closure()\\

Boolean reachability matrix: reach[i][j] is true iff j is reachable from i.

\textbf{Why:} You need to know if vertex j is reachable from vertex i. \par \textbf{Where to use:} Use for reachability queries, dependency analysis, and type hierarchy checking. Boolean matrix.

\begin{lstlisting}
dsa::graph<int> g(3);
g.add_edge(0, 1); g.add_edge(1, 2);
auto tc = g.transitive_closure();
// tc[0][2] = true
\end{lstlisting}

\subsubsection{DAG Shortest Path}
\textbf{Function:} g.dag\_shortest\_path(start)\\

Processes vertices in topological order for optimal single-pass shortest path.

\textbf{Why:} You need shortest paths in a DAG efficiently. \par \textbf{Where to use:} Use for project scheduling, critical-path analysis, and DAG optimization. It processes vertices once in topological order, in \(O(V+E)\) time.

\begin{lstlisting}
dsa::graph<int> g(4);
g.add_edge(0, 1, 1); g.add_edge(0, 2, 4);
g.add_edge(1, 3, 2); g.add_edge(2, 3, 1);
auto dist = g.dag_shortest_path(0);
// dist = {0, 1, 4, 3}
\end{lstlisting}

\subsubsection{Eulerian Cycle}
\textbf{Function:} g.eulerian\_cycle()\\

Hierholzer's algorithm for finding a cycle that uses every edge exactly once. Requires in-degree = out-degree for all vertices.

\textbf{Why:} You need a cycle that uses every edge exactly once. \par \textbf{Where to use:} Use for genome assembly, Chinese postman problem, and graph drawing. Requires in-degree = out-degree for all vertices.

\begin{lstlisting}
dsa::graph<int> g(3);
g.add_edge(0, 1); g.add_edge(1, 2); g.add_edge(2, 0);
auto cycle = g.eulerian_cycle();
// cycle = {0, 1, 2, 0} (or rotation)
\end{lstlisting}

% ---- MST ----
\subsection{Minimum Spanning Tree}

\subsubsection{Kruskal's MST}
\textbf{Header:} code/include/graph.h\\
\textbf{Function:} dsa::kruskal\_mst(n, edges)\\

Sorts edges by weight, greedily adds cheapest edges that don't form cycles (using Union-Find).

\textbf{Why:} You need the minimum weight spanning tree using greedy edge selection. \par \textbf{Where to use:} Use for network design, clustering, and Steiner tree approximation. Sort edges greedily with Union-Find for cycle detection.

\begin{lstlisting}
std::vector<dsa::weighted_edge> edges =
 {{0,1,4},{1,2,2},{0,2,3}};
auto mst = dsa::kruskal_mst(3, edges);
// mst has 2 edges, total weight = 5
\end{lstlisting}

\subsubsection{Boruvka's MST}
\textbf{Header:} code/include/graph.h\\
\textbf{Function:} dsa::boruvka\_mst(n, edges)\\

Each component picks its cheapest outgoing edge. Parallelizable.

\textbf{Why:} You need a parallelizable MST algorithm. \par \textbf{Where to use:} Use for parallel MST computation. Each component picks cheapest outgoing edge. Parallelizable unlike Prim's.

\begin{lstlisting}
std::vector<dsa::weighted_edge> edges =
 {{0,1,4},{1,2,2},{0,2,3}};
auto mst = dsa::boruvka_mst(3, edges);
// mst has 2 edges, total weight = 5
\end{lstlisting}

\subsubsection{Prim's MST}
\textbf{Header:} code/include/prim.h\\
\textbf{Class:} dsa::prim\_mst$<$...$>$\\

Grows a single tree from a source vertex by always adding the cheapest edge connecting the tree to a non-tree vertex.

\textbf{Why:} You need MST growing from a single source vertex. \par \textbf{Where to use:} Use for network design, cluster analysis. Grows tree by always adding cheapest edge to non-tree vertex.

\begin{lstlisting}
dsa::prim_mst pm(4);
pm.add_edge(0, 1, 4); pm.add_edge(1, 2, 2);
pm.add_edge(2, 3, 3); pm.add_edge(0, 3, 5);
auto res = pm.compute();
// res.total_weight = 9
\end{lstlisting}

% ---- SCC ----
\subsection{Strongly Connected Components}
\textbf{Header:} code/include/scc.h

\subsubsection{Kosaraju's Algorithm}
\textbf{Class:} dsa::kosaraju\_scc\\

Two-pass DFS: first pass gets finish order on original graph, second pass finds components on transpose.

\textbf{Why:} You need to find strongly connected components in a directed graph. \par \textbf{Where to use:} Use for dependency analysis, web page clustering, and graph simplification. Two-pass DFS on original and transpose.

\begin{lstlisting}
dsa::kosaraju_scc ks(5);
ks.add_edge(0, 1); ks.add_edge(1, 2);
ks.add_edge(2, 0); ks.add_edge(2, 3);
ks.add_edge(3, 4);
auto sccs = ks.find_sccs();
// sccs = {{0,1,2}, {3}, {4}}
\end{lstlisting}

\subsubsection{Tarjan's Algorithm}
\textbf{Class:} dsa::tarjan\_scc\\

Single-pass DFS using a stack. Finds SCCs by tracking discovery times and low-link values.

\textbf{Why:} You need SCCs with a single DFS pass (no transpose graph needed). \par \textbf{Where to use:} Use for SCCs when you want single-pass simplicity. Uses discovery times and low-link values.

\begin{lstlisting}
dsa::tarjan_scc ts(5);
ts.add_edge(0, 1); ts.add_edge(1, 2);
ts.add_edge(2, 0); ts.add_edge(2, 3);
ts.add_edge(3, 4);
auto sccs = ts.find_sccs();
// sccs = {{4}, {3}, {0,1,2}}
\end{lstlisting}

% ---- MAX FLOW ----
\subsection{Maximum Flow \& Matching}

\subsubsection{Dinic's Algorithm}
\textbf{Header:} code/include/dinic.h\\
\textbf{Class:} dsa::dinic$<$FlowUnit$>$\\

Builds level graphs via BFS, then finds blocking flows via DFS. Much faster in practice than Ford-Fulkerson on unit-capacity networks.

\textbf{Why:} You need maximum flow with good practical performance. \par \textbf{Where to use:} Use as the default max-flow algorithm. BFS level graphs + DFS blocking flows. Much faster than Ford-Fulkerson in practice.

\begin{lstlisting}
dsa::dinic<int> d(4);
d.add_edge(0, 1, 10);
d.add_edge(0, 2, 10);
d.add_edge(1, 3, 10);
d.add_edge(2, 3, 10);
int flow = d.max_flow(0, 3);
// flow = 20
\end{lstlisting}

\subsubsection{Ford-Fulkerson}
\textbf{Header:} code/include/flow\_matching.h\\
\textbf{Class:} dsa::ford\_fulkerson(n)\\

Augments along augmenting paths until no more exist. Complexity depends on max flow value.

\textbf{Why:} You want the simplest augmenting-path max-flow algorithm. \par \textbf{Where to use:} Use for understanding max-flow concepts. Its running time depends on path selection and capacity values; use Dinic's or Edmonds--Karp for stronger bounds.

\begin{lstlisting}
dsa::ford_fulkerson ff(4);
ff.add_edge(0, 1, 10); ff.add_edge(0, 2, 10);
ff.add_edge(1, 3, 10); ff.add_edge(2, 3, 10);
int flow = ff.max_flow(0, 3);
// flow = 20
\end{lstlisting}

\subsubsection{Edmonds-Karp}
\textbf{Header:} code/include/flow\_matching.h\\
\textbf{Class:} dsa::edmonds\_karp(n)\\

Ford-Fulkerson using BFS for shortest augmenting path. Guaranteed polynomial time.

\textbf{Why:} You want Ford--Fulkerson with a guaranteed polynomial-time bound. \par \textbf{Where to use:} Use when BFS shortest augmenting paths are sufficient and the graph is small. It is simpler than Dinic's algorithm but can be slower.

\begin{lstlisting}
dsa::edmonds_karp ek(4);
ek.add_edge(0, 1, 10); ek.add_edge(0, 2, 10);
ek.add_edge(1, 3, 10); ek.add_edge(2, 3, 10);
int flow = ek.max_flow(0, 3);
// flow = 20
\end{lstlisting}

\subsubsection{Hopcroft-Karp (Bipartite Matching)}
\textbf{Header:} code/include/flow\_matching.h\\
\textbf{Class:} dsa::bipartite\_matching(n\_left, n\_right)\\

Finds maximum matching in bipartite graphs by finding multiple augmenting paths per BFS phase.

\textbf{Why:} You need maximum matching in a bipartite graph. \par \textbf{Where to use:} Use for job assignment, resource allocation, and matching problems. Hopcroft--Karp improves augmenting-path search by finding multiple paths per BFS phase.

\begin{lstlisting}
dsa::bipartite_matching bm(3, 3);
bm.add_edge(0, 0);
bm.add_edge(0, 1);
bm.add_edge(1, 1);
bm.add_edge(2, 2);
int sz = bm.max_matching();
// sz = 3
\end{lstlisting}

\subsubsection{Hungarian Algorithm}
\textbf{Header:} code/include/hungarian.h\\
\textbf{Class/Function:} dsa::hungarian, dsa::hungarian\_min(cost\_matrix)\\

Finds minimum-cost perfect matching in a bipartite graph (assignment problem).

\textbf{Why:} You need minimum-cost perfect matching in a bipartite graph. \par \textbf{Where to use:} Use for assignment problems such as workers to tasks, exam seating, and ride-sharing. The algorithm returns an optimal assignment for its supported square cost matrix.

\begin{lstlisting}
std::vector<std::vector<int>> cost =
 {{1, 2, 3}, {4, 5, 6}, {7, 8, 9}};
auto res = dsa::hungarian_min(cost);
// res.cost = 15, res.assignment = {0,1,2}
\end{lstlisting}

%% ============================================================
%% ADDITIONAL ALGORITHMS
%% ============================================================

\section{Additional Algorithms}

\subsubsection{de Bruijn Graph \& Sequence}
\textbf{Header:} code/include/de\_bruijn.h\\
\textbf{Function:} dsa::de\_bruijn\_sequence(k, n)\\

Generates a cyclic sequence where every possible substring of length n over an alphabet of size k appears exactly once.

\textbf{Why:} You need a cyclic sequence containing every possible substring of length n. \par \textbf{Where to use:} Use for genome assembly, lock combinations, and pseudo-random number generation. Every n-length substring appears exactly once.

\begin{lstlisting}
auto seq = dsa::de_bruijn_sequence(2, 3);
// seq contains all 3-bit binary substrings
\end{lstlisting}

\subsubsection{Landau-Vishkin}
\textbf{Header:} code/include/additional\_algorithms.h\\
\textbf{Function:} dsa::landau\_vishkin(a, b, k)\\

Computes edit distance only up to threshold k. Returns early if distance exceeds k.

\textbf{Why:} You need edit distance only up to threshold k, and want early termination. \par \textbf{Where to use:} Use for spelling correction, near-duplicate detection, and DNA comparison when you only care about small differences.

\begin{lstlisting}
auto matches = dsa::landau_vishkin(
 "hello world", "helo word", 3);
// matches = {positions within edit distance 3}
\end{lstlisting}

\subsubsection{Hirschberg LCS}
\textbf{Header:} code/include/additional\_algorithms.h\\
\textbf{Function:} dsa::hirschberg\_lcs(a, b)\\

Space-efficient LCS computation. Divides and conquers using linear-space needleman-wunsch.

\textbf{Why:} You need LCS using linear auxiliary space instead of a full \(O(mn)\) table. \par \textbf{Where to use:} Use for sequence alignment on memory-constrained devices, diff algorithms, and bioinformatics where space is limited.

\begin{lstlisting}
auto lcs = dsa::hirschberg_lcs("ABCBDAB", "BDCAB");
// lcs = "BCAB" (O(min(m,n)) space)
\end{lstlisting}

\subsubsection{Shortest Common Superstring}
\textbf{Header:} code/include/additional\_algorithms.h\\
\textbf{Function:} dsa::shortest\_common\_superstring(strings)\\

Uses a greedy overlap heuristic to construct a short string containing all input strings as substrings; it is not guaranteed to be globally shortest.

\textbf{Why:} You need a short common superstring and can accept an approximation. \par \textbf{Where to use:} Use for educational genome-assembly and text-reconstruction experiments; use an exact or domain-specific method when optimality is required.

\begin{lstlisting}
auto scs = dsa::shortest_common_superstring(
 {"ABC", "BCA", "CAB"});
// scs is a short common superstring, not necessarily the shortest
\end{lstlisting}

\subsubsection{Crochemore-Perrin Two-Way}
\textbf{Header:} code/include/additional\_algorithms.h\\
\textbf{Function:} dsa::crochemore\_perrin(text, pattern)\\

Optimal two-way pattern matching using critical factorization. Text is scanned left-to-right.

\textbf{Why:} You need optimal two-way pattern matching. \par \textbf{Where to use:} Use for optimal string matching with left-to-right text scan. Critical factorization enables with minimal preprocessing.

\begin{lstlisting}
auto pos = dsa::crochemore_perrin(
 "ABCABCABC", "ABCABC");
// pos = {0, 3}
\end{lstlisting}

%% ============================================================
%% QUICK REFERENCE TABLE
%% ============================================================

\section{New Industrially Relevant Implementations}

The following additions connect the guide's algorithms to reusable, tested C++23 headers. They are deliberately small canonical APIs that can be adapted to production systems.

\begin{itemize}
\item \texttt{a\_star.h}: \texttt{dsa::a\_star(graph, start, goal, heuristic)} returns an optional distance and reconstructed path. Use non-negative edge weights and an admissible heuristic for optimality.
\item \texttt{graph.h}: \texttt{articulation\_points()} and \texttt{bridges()} use Tarjan low-link analysis on undirected simple graphs, useful for network-dependency and failure analysis.
\item \texttt{min\_cost\_flow.h}: \texttt{dsa::min\_cost\_max\_flow} supports capacities, signed edge costs, flow limits, and returns both total flow and cost.
\item \texttt{persistent\_segment\_tree.h}: \texttt{update(version, index, value)} creates a new version while preserving older versions; \texttt{query(version, left, right)} uses half-open ranges.
\item \texttt{cuckoo\_filter.h}: \texttt{dsa::cuckoo\_filter} provides expected constant-time approximate membership checks and deletion. It may return false positives, never false negatives after a successful insertion, and insertion can fail when full.
\item \texttt{heavy\_hitters.h}: \texttt{dsa::misra\_gries} tracks candidate frequent keys in bounded space; candidates require a second pass or exact counter for final frequency guarantees.
\item \texttt{wavelet\_tree.h}: \texttt{dsa::wavelet\_tree} supports byte-string access, rank, select, and range counting, useful for indexed text and compressed analytics.
\end{itemize}

\section{Quick Reference}

\begin{center}
\small
\begin{tabular}{@{}lll@{}}
\toprule
\textbf{Header} & \textbf{Category} & \textbf{Complexity} \\
\midrule
sorting.h & Sorting (12 algos) & varies by algorithm \\
graph.h & Graph (15+ algos) & varies by algorithm \\
a\_star.h & A* search & typical $O((V+E)\log V)$ \\
min\_cost\_flow.h & Min-cost max-flow & successive shortest paths \\
dynamic\_programming.h & DP (4 algos) & varies by algorithm \\
greedy.h & Greedy (2 algos) & varies by algorithm \\
divide\_conquer.h & D\&C (2 algos) & varies by algorithm \\
backtracking.h & Backtracking (3) & exponential \\
exact\_matching.h & String (5 algos) & linear / sublinear \\
approximate\_matching.h & String (4) & varies by algorithm \\
multiple\_matching.h & String (2) & varies by algorithm \\
string\_compression.h & Compression (4) & varies by algorithm \\
dinic.h & Max Flow & algorithm-dependent \\
flow\_matching.h & Flow/Match (3) & algorithm-dependent \\
probabilistic.h & Probabilistic (3) & sublinear space \\
bloom\_filter.h & Probabilistic & expected constant \\
cuckoo\_filter.h & Probabilistic membership & expected constant \\
heavy\_hitters.h & Streaming candidates & (O(1)) amortized update \\
avl\_tree.h & Balanced BST & \(O(\log n)\) \\
red\_black\_tree.h & Balanced BST & \(O(\log n)\) \\
b\_tree.h & Multi-way tree & \(O(\log n)\) \\
fenwick\_segment.h & Range query & \(O(\log n)\) \\
segment\_tree\_lazy.h & Range query & \(O(\log n)\) \\
persistent\_segment\_tree.h & Versioned range query & \(O(\log n)\) update/query \\
union\_find.h & Disjoint set & \(O(\alpha(n))\) amortized \\
kd\_tree.h & Spatial & varies by query \\
bwt.h & BWT/FM-index & varies by query \\
wavelet\_tree.h & Indexed text & \(O(\log \sigma)\) rank \\
nfa.h & Regex & linear in text length \\
\bottomrule
\end{tabular}
\end{center}

\vfill
\begin{center}
\small\textit{All implementations live in namespace dsa, use C++23 features,
and are header-only.}\\
\textit{Build: cd code \&\& cmake -B build \&\& cmake --build build \&\& ctest --test-dir build}
\end{center}

\end{document}
